%%%%%%%%%%%%%%%%%%%%%%%%%%%%%%%%%%%%%%%%%%%%%%%%%%%%%%%%%%%%%%%%%%%%%%%%%%%%%%%%

%-------------------------------------------------------------------------------
%       PACKAGES AND OTHER DOCUMENT CONFIGURATIONS
%-------------------------------------------------------------------------------

%%%%%%%%%%%%%%%%%%%%%%%%%%%%%%%%%%%%%%%%%%%%%%%%%%%%%%%%%%%%%%%%%%%%%%%%%%%%%%%%
% You can have multiple style options the legal options ones are:
%
%   centered:   the name and address are centered at the top of the page 
%                               (default)
%
%   line:               the name is the left with a horizontal line then the address to
%                               the right
%
%   overlapped: the section titles overlap the body text (default)
%
%   margin:             the section titles are to the left of the body text
%
%   11pt:               use 11 point fonts instead of 10 point fonts
%
%   12pt:               use 12 point fonts instead of 10 point fonts
%
%%%%%%%%%%%%%%%%%%%%%%%%%%%%%%%%%%%%%%%%%%%%%%%%%%%%%%%%%%%%%%%%%%%%%%%%%%%%%%%%

\documentclass[margin]{res}  

% Default font is the helvetica postscript font
\usepackage{helvet}

% Increase text height
\textheight=700pt

\begin{document}

%-------------------------------------------------------------------------------
%       NAME AND ADDRESS SECTION
%-------------------------------------------------------------------------------
\name{Victor Delaplaine}

% Note that addresses can be used for other contact information:
% -phone numbers
% -email addresses
% -linked-in profile

\address{
\\Website : web.calpoly.edu/~vdelapla/
\\LinkedIn : www.linkedin.com/in/victor-delaplaine
\\Github : www.github.com/ByVictorrr\\
}
\address{vdelapla@calpoly.edu \\(650) 740-6691\\}

% Uncomment to add a third address
%\address{Address 3 line 1\\Address 3 line 2\\Address 3 line 3}
%-------------------------------------------------------------------------------

\begin{resume}

%-------------------------------------------------------------------------------
%       EDUCATION SECTION
%-------------------------------------------------------------------------------
\section{OBJECTIVE}
{\sl Senior level Computer Engineering Student at Cal Poly. Inquisitive, hard-working and consistent. Looking for internship opportunities where I can apply my skills and contribute to real-world projects.}

\section{EDUCATION}

\textbf{California Polytechnic State University}, San Luis Obispo\\
{\sl Bachelor of Science (BS)}, Computer Engineering(CPE)\\
Expected June, 2020
\hfill Overall GPA: 3.897/4.00 


%-------------------------------------------------------------------------------
%       COMPUTER SKILLS SECTION
%-------------------------------------------------------------------------------
\section{TECHNICAL\\SKILLS}

\textbf{Programming:} Shell, C++, C, Java, Verilog, SystemVerilog
\\
\textbf{Markdown :} MarkDown, LaTex, HTML
\\
\textbf{Familiar :} Matlab, CSS, MySQL, Regular Expressons
\\
\textbf{General :} Data Structures, OOP and Assembly Programming, Digital Design, Computer Design, git, UNIX

%-------------------------------------------------------------------------------
% Modify the format of each position
\begin{format}
\title{l}\\
\dates{l}\location{r}\\
\body\\
\end{format}
%-------------------------------------------------------------------------------




%-------------------------------------------------------------------------------
%      EXPERIENCE
%-------------------------------------------------------------------------------
\section{EXPERIENCE}

\textbf{San Francisco State : Machine learning lead intern \hfill{June 18 - August 18}\\} \\
\normalfont{ 
-\ Assisted in developing an android application, for a Neural Machine Interface (NMI)
\\
-\ Set up an AWS server to get a working SQL base in order to store users data
\begin{itemize}
\item \textbf{Tools :} MySQL, teamwork, AWS, research and presenting
\end{itemize}
\begin{itemize}
\item \textbf{Link :} github.com/alexanderadavid/Myo-HMI-Android
\end{itemize}

\textbf{Canada College : STEM tutor \hfill{January 18 - June 18}\\}
\\
\normalfont{
    - Math, Physics, Engineering, and Computer Science tutor \\
    - Enjoyed helping students succeed by showing them how to solve problems 

}
\begin{itemize}
\item \textbf{Tools :} Teaching
\end{itemize}
\begin{itemize}
\item \textbf{Link :} https://www.canadacollege.edu/stemcenter/ 
\end{itemize}



\section{RECENT PROJECTS}

% Project 1
\employer{ByVictorrr}
\location{}
\title{\textbf{aProjector - Autofocusing Projector \hfill September 2018}}
\begin{position}
% \begin{itemize}
 The aProjector, was a invented by a classmate and I. It is simply a projector which projects a phone’s screen on the wall and autofocuses it. I referred to the thin lens equation as a basis to develop a program written in C++ run on the Arduino Mega 2560. This Arduino controls the motor to turn the stepper motor into a given angle needed to produce a clear image on the wall.
% \end{itemize}https://www.overleaf.com/project/5c3cda8d86f90560469d48be
\begin{itemize}
\item \textbf{Technology/Tools:} C/C++, LaTex, markdown
\end{itemize}
\begin{itemize}
\item \textbf{Link :} github.com/byvictorrr/aProjector
\end{itemize}
\end{position}




% Project 2
\employer{ByVictorrr}
\location{}
\title{\textbf{ SolarRAT - Solar Tracker that is RAT fast  \hfill March 2019}}
\begin{position}
% \begin{itemize}
 This is a solar tracker I designed using an arduino and a RAT MCU. The solar tracker sweeps through 12 position via servo motor and records the intensity of light at each of the 12 locations. Using a bubble sort, the highest intensity of light location is found and the the solar panel stays there. Below shows a video in the link.
% \end{itemize}https://www.overleaf.com/project/5c3cda8d86f90560469d48be
\begin{itemize}
\item \textbf{Technology/Tools:} C++, Arduino, Computer Design, Assembly Language, SystemVerilog, Flowcharts
\end{itemize}
\begin{itemize}
\item \textbf{Link :} github.com/byvictorrr/SolarRAT
\end{itemize}
\end{position}

% Project 3
\employer{ByVictorrr}
\location{}
\title{\textbf{Assembler - Mips assembler \hfill April 2019}}
\begin{position}
% \begin{itemize}
Designed an assembler that transforms a large set of instructions into the mips machine code.
% \end{itemize}https://www.overleaf.com/project/5c3cda8d86f90560469d48be
\begin{itemize}
\item \textbf{Technology/Tools:} Java
\end{itemize}
\begin{itemize}
\item \textbf{Link :} https://github.com/ByVictorrr/mips-assembler
\end{itemize}
\end{position}

% Project 4
\employer{ByVictorrr}
\location{}
\title{\textbf{ RAT MCU - Designed a basic computer  \hfill March 2019}}
\begin{position}
% \begin{itemize}
I designed a fully functional computer that has around 40 instruction. This computer consists of basic modules - Program Rom, Register file, RAM, Program Counter, Stack Pointer, Arithmetic and Logic Unit, Flags, Shadow Flags, Interrupt Flag and a Control Unit. These were all modeled using SystemVerilog on a Basys3 FPGA board. 
% \end{itemize}https://www.overleaf.com/project/5c3cda8d86f90560469d48be
\begin{itemize}
\item \textbf{Technology/Tools:} Digital Design, Computer Design, Assembly Language, SystemVerilog
\end{itemize}
\begin{itemize}
\item \textbf{Link :} github.com/byvictorrr/CPE233
\end{itemize}
\end{position}




\section{CERTIFICATION}
\par
\normalfont{}
\\
\normalfont{}\textbullet{} \textbf{Computer Science - C++ Certificate of Achievement} by Canada College
\\ 
\textbullet{} \textbf{Engineering - Certificate of Achievement} by Canada College 


\section{RELEVANT\\COURSES}
\par

\normalfont{\textbullet{} PCB CAD: github.com/byvictorrr/IME156 }
\\
\normalfont{\textbullet{} Digital Design: github.com/byvictorrr/CPE133 }
\\
\normalfont{\textbullet{} Object Orientated Programming in Java: github.com/byvictorrr/CPE203 }
\\
\normalfont{\textbullet{} Computer Design and Assembly Language Programming: github.com/byvictorrr/CPE233 }
\normalfont{\textbullet{} Computer Architecture: github.com/byvictorrr/CPE315 }
\\
\normalfont{\textbullet{} Systems Programming: github.com/byvictorrr/CPE357 }
\\
\normalfont{\textbullet{} Object Orientated Programming in C++: https://github.com/ByVictorrr/Sudoku }
\normalfont{\textbullet{} Data Structures in C++: https://github.com/ByVictorrr/GhostGame}
\\
\normalfont{\textbullet{} Semiconductor Device Electronics: https://github.com/ByVictorrr/EE306}
\\
\normalfont{\textbullet{} Digital Electronics and Integrated Circuits : https://github.com/ByVictorrr/EE307}


\end{resume}
\(\)\end{document}
